\section{Gobernanza de datos: derechos de autor, privacidad y control de calidad}

En cuanto la industria de datos alcanza escala, los problemas de gobernanza se vuelven ineludibles. Los datos para LLM implican derechos de autor, privacidad y transparencia sobre su procedencia; los datos para robots implican la recolección en entornos reales, los escenarios de los usuarios, los registros de sensores y la responsabilidad en materia de seguridad; los datos de simulación, por su parte, implican supuestos físicos y sesgos de evaluación. Sin gobernanza, cuantos más datos haya, mayor será el riesgo.

\subsection{El triángulo de la gobernanza}

\noindent\begin{longtable}{p{0.22\textwidth}p{0.34\textwidth}p{0.35\textwidth}}
\toprule
Dimensión & Pregunta que debe responderse & Consecuencia de fallar \\
\midrule
Derechos de autor y licencias & ¿Los datos tienen una procedencia legal y un alcance de uso definido? & Controversias sobre los datos de entrenamiento y riesgos comerciales posteriores. \\
Privacidad y seguridad & ¿Los datos contienen información personal, ubicaciones, hogares, fábricas o secretos comerciales? & Filtración de información sensible e imposibilidad de entrar en entornos empresariales. \\
Calidad y trazabilidad & ¿De dónde provienen los datos, cómo se limpian y cómo influyen en el modelo? & No es posible localizar las fallas del modelo y la Recipe no es reproducible. \\
\bottomrule
\end{longtable}

\begin{warningbox}{La gobernanza de datos no es un trámite legal de último momento}
Si la procedencia, la limpieza y las licencias de los datos no son trazables, cuanto más capaz sea el modelo, mayor será el riesgo. Esto ocurre sobre todo en los robots y en los entornos empresariales, porque los datos suelen contener espacios reales, equipos reales y comportamientos reales de los usuarios.
\end{warningbox}

\subsection{Las cuatro preguntas del control de calidad}

Antes de que un conjunto de datos entre en el entrenamiento, hay que hacerse al menos cuatro preguntas: si representa la tarea objetivo; si contiene las capacidades que le faltan al modelo; si tiene etiquetas o retroalimentación confiables; si introducirá atajos erróneos. En el caso de los datos de simulación, además hay que preguntar: ¿la política aprendida en la simulación puede transferirse al mundo real?

\begin{importantbox}{El control de calidad es requisito previo de la Recipe}
Sin control de calidad, la Recipe no es más que una proporción de mezcla elegida casi por superstición; con control de calidad, la Recipe puede convertirse en un método de ingeniería reproducible.
\end{importantbox}

\subsection{Resumen de la sección}

La gobernanza de datos determina si los datos pueden usarse a largo plazo. Los derechos de autor, la privacidad, la trazabilidad y el control de calidad son el umbral que la industria de datos debe cruzar para pasar de la «recolección» a la «infraestructura».
